\section{Explicit bounds for resource-coupled followers}
\label{app:quantitative}

The estimates below give constants of polynomial bit length, though not
necessarily of small numerical value, for the accuracy algorithms. They
instantiate classical polyhedral error bounds \cite{Hoffman1952}; no optimal
Hoffman constant is needed.

\begin{lemma}[Exact resource projection]\label{lem:resource-projection}
For a rational matrix $C\in\Q^{k\times N}$, with fixed $k$, the set
\[
 X_F=\{x\in X:\exists z\in[0,1]^N,\ Cz\le b(x)\}
\]
has a polynomial-size rational inequality description when $X$ is a
rational polytope and $b$ is rational affine. That description and its
feasibility can be computed in polynomial bit time.
\end{lemma}
\begin{proof}
The right-hand sides that admit a feasible follower form the closed
polyhedron $C[0,1]^N+\R_+^k$. By separation, membership is equivalent to
\begin{equation}\label{eq:resource-support}
 \lambda\trans b(x)\ge\sum_{i=1}^N\min\{0,\lambda\trans C_i\}
                              \qquad\text{for all }\lambda\ge0.
\end{equation}
Indeed, a finite lower support functional on the orthant sum must be
nonnegative, and its minimum over the box is the displayed sum. This sum is
linear on every cone of the central arrangement $\lambda\trans C_i=0$ in
the nonnegative orthant. Such cones are pointed, so it suffices to check
their extreme rays. Each ray has $k-1$ independent active hyperplanes among
the column and coordinate hyperplanes. Enumerate all such subsets, compute
their one-dimensional nullspaces, and retain each nonnegative direction.
The resulting $O((N+k)^{k-1})$ tests include all necessary rays; extra
nonnegative directions give valid inequalities. Cramer's rule supplies polynomial
rational bit lengths. Coordinate hyperplanes cover rank deficiency and
zero columns, including lower-dimensional cones. For $k=1$, the empty
subset gives the ray $+1$; for $k=0$, $X_F=X$. Substitute $b(x)$
in these finitely many inequalities and use rational LP.
\end{proof}

\begin{lemma}[Uniform multipliers and feasibility repair]
\label{lem:resource-constants}
Assume $N\ge1$ and put
\begin{equation}\label{eq:integer-repair-constants}
 \begin{gathered}
 \mathsf A=d_0[C;I;-I],\qquad
 M=\max\{1,\max_{ij}|\mathsf A_{ij}|\},\\
 V=N!(1+NM^2)^N,\qquad K=d_0N^3MV,
 \end{gathered}
\end{equation}
where the positive integer $d_0$ clears all denominators of $C$.
These constants have polynomial bit length. For every $x\in X_F$
and $q\in[0,1]^N$ with $Cq-b(x)\le\delta\mathbf1$, $\delta\ge0$,
there is a feasible $y$ with
\begin{equation}\label{eq:resource-repair}
                         \norm{q-y}_\infty\le K\delta.
\end{equation}
If a differentiable convex follower objective has gradient bounded by
$\overline G$ in infinity norm on the boxes, each optimum has resource
KKT multipliers in $[0,\Lambda]^k$, where
$\Lambda=\max\{1,K\overline G\}$. No strict feasibility is required.
\end{lemma}
\begin{proof}
By the conic support reduction used in \cref{lem:local-branches}, a
vector $v$ in a cone generated by rows of the integer matrix $\mathsf A$
has a nonnegative representation using linearly independent rows
$\mathsf A_J$, with $|J|\le N$. Its coefficients are
\[
 \alpha=(\mathsf A_J\mathsf A_J\trans)^{-1}\mathsf A_Jv.
\]
The Gram determinant is a positive integer. Its entries have magnitude at
most $NM^2$, and by the permutation expansion every cofactor, including the
zero-order cofactor one, has magnitude at most $V$. Hence
\[
 \norm{\alpha}_1\le N^3MV\norm{v}_2.
\]
For $\norm{v}_\infty\le\overline G$, the same formula bounds each
individual coefficient by $N^2MV\overline G\le N^3MV\overline G$.
At an optimum apply this to $-\nabla F_x(z^*)$, using only active
normals. Undoing their scale multiplies coefficients by $d_0$, so
resource multipliers are at most $K\overline G$. The active-normal
formula is valid for every polyhedron, including degenerate faces.
Convexity makes the resulting KKT conditions sufficient.

For repair, let $y$ be the Euclidean projection of $q$ onto the
nonempty follower polytope and apply the representation to $v=q-y$, which
lies in the cone of the scaled rows active at $y$. For these rows,
$\mathsf A_j(q-y)$ is the violation at $q$. Resource violations are at
most $d_0\delta$, and box violations are nonpositive. Thus
\[
 \norm{v}_2^2=\alpha\trans\mathsf A_J(q-y)
 \le d_0\delta\norm{\alpha}_1\le K\delta\norm{v}_2.
\]
For $v\ne0$ divide by $\norm{v}_2$; for $v=0$ the result is immediate.
The bit bounds follow from $\log N!=O(N\log N)$ and the explicit
integer powers. The accuracy algorithm computes neither the projection nor
the correct active set.
\end{proof}

\begin{lemma}[Polynomial growth without strong curvature]
\label{lem:polynomial-bregman}
Let $g_i(0)=0$, $g_i(1)=G_i>0$ be strictly increasing rational
polynomials of degrees at most $P\ge1$, and $f_i'=g_i$. For the
convex aggregate model \eqref{eq:accuracy-follower}, at every fixed
follower-feasible leader and all $y,z\in[0,1]^N$,
\begin{equation}\label{eq:accuracy-bregman}
 F_x(y)-F_x(z)-\nabla F_x(z)\trans(y-z)
                 \ge\mu\norm{y-z}_\infty^{P+1},\qquad
 \mu=\frac{\min_iG_i}{4(4P)^P}.
\end{equation}
If every coefficient of each $g_i$ is nonnegative, one may instead
use the sharper $\mu=\min_iG_i/(P+1)$.
\end{lemma}
\begin{proof}
For one marginal of actual degree $d$, let $h=|y_i-z_i|$.
In either direction, by \cref{lem:inverse-modulus} on a response
interval of length $h/2$, the Bregman integral contains an interval of
length $h/2$ on which the marginal differs from its base value by at least
$G_i(h/(4d))^d/2$. Its integral is therefore at least
$G_ih^{d+1}/[4(4d)^d]\ge G_ih^{P+1}/[4(4P)^P]$.
For nonnegative coefficients, the Bregman divergence of
$t^{j+1}/(j+1)$ is at least $|y_i-z_i|^{j+1}/(j+1)$: in the
forward direction integrate $(z_i+t)^j-z_i^j\ge t^j$; in the
reverse direction integrate $z_i^j-(z_i-t)^j\ge t^j$.
Both follow from the binomial theorem with nonnegative arguments.
Multiply by the coefficients and use $h\le1$ and $j\le P$.
Sum over coordinates; the affine terms cancel, and the aggregate has
nonnegative Bregman divergence by convexity. The coordinate of largest
displacement gives the asserted infinity-norm bound.
\end{proof}

Polynomial solution-map regularity has a direct predecessor in
Jeyakumar, Lasserre, Li, and Ph\d{a}m~\cite[Theorem 2.3]{JeyakumarEtAl2016}.
For a compact polynomial follower feasible set independent of the leader,
their theorem gives a one-sided inclusion of the nearby solution set
$Y(x)$ in a H\"older neighborhood of $Y(\bar x)$ at a fixed reference
leader $\bar x$, with an exponent bound depending on polynomial degree,
follower dimension and constraint count. A closer precedent for the exponent
is Wu et al.~\cite[Lemma 4.2 and Appendix B.1]{WuEtAl2026}: for an
unconstrained follower with uniform convexity of order $p$ and a lower
gradient that is Lipschitz in the leader, they obtain the pairwise exponent
$1/(p-1)$, so $p=P+1$ already yields $1/P$ in that setting. For the
structured unique-response class here, the following proposition keeps this
exponent when affine resource right-hand sides vary with the leader and gives
a computable rational constant of polynomial encoding length. The exponent
is independent of follower dimension and aggregate degree.

\begin{proposition}[Leader-response modulus and a sharper exponent]
\label{prop:accuracy-response-modulus}
For \eqref{eq:accuracy-follower} with $N\ge1$, effective rational constants with
polynomial encoding give, for $x,x'\in X_F$ and
$\Delta=\norm{x-x'}_\infty$,
\begin{equation}\label{eq:response-modulus-basic}
 \norm{z(x)-z(x')}_\infty
 \le KB_b\Delta+
       [2(T_x+L_FKB_b)\Delta/\mu]^{1/(P+1)}.
\end{equation}
Here $B_b$ bounds the Lipschitz constant of $b$ and
$|F_x(z)-F_{x'}(z)|\le T_x\Delta$ uniformly on the box.
If $\phi=0$, replace $2(T_x+L_FKB_b)$ by
$2L_FKB_b+B_\ell$, where
$B_\ell=\sum_i\norm{\nabla\ell_i}_1$.
In fact there is an effective rational constant $C_z$ of polynomial
encoding such that
\begin{equation}\label{eq:response-modulus-sharp}
                  \norm{z(x)-z(x')}_\infty\le C_z\Delta^{1/P}.
\end{equation}
The exponent $1/P$ is best possible for this class.
\end{proposition}
\begin{proof}
For the first bound, repair $z(x)$ to $y$ feasible at $x'$, and
$z(x')$ to $y'$ feasible at $x$, each within $KB_b\Delta$.
Optimality and the two Lipschitz bounds give
\[
 \begin{split}
 F_{x'}(y)&\le F_x(z(x))+T_x\Delta+L_FKB_b\Delta\\
 &\le F_x(y')+T_x\Delta+L_FKB_b\Delta\\
 &\le F_{x'}(z(x'))+2(T_x+L_FKB_b)\Delta.
 \end{split}
\]
Apply \cref{lem:polynomial-bregman} at the constrained optimum
$z(x')$ and add the repair distance. For $\phi=0$, retaining the
cancellation of the leader-linear terms bounds the gap instead by
\[
 2L_FKB_b\Delta+
 |(\ell(x)-\ell(x'))\trans(z(x)-z(x'))|
                         \le(2L_FKB_b+B_\ell)\Delta.
\]
These estimates need only coefficient sums in the fixed aggregate
variables and in the local marginals.

The sharper assertion needs a separate argument. Let $A_x,A_z\ge0$ be
rational bounds such that the full follower gradient obeys
\[
 \norm{\nabla F_x(z)-\nabla F_{x'}(z)}_\infty\le A_x\Delta,
 \qquad
 \norm{\nabla F_x(z)-\nabla F_x(z')}_\infty\le A_z\norm{z-z'}_\infty.
\]
Such bounds of polynomial encoding length can be computed without
expanding in $N$ variables: local derivative coefficient sums, bounds on
the aggregate Hessian and its mixed leader derivatives, and the row and
column sums of $U$ suffice. Set $D_b=KB_b$.

First suppose $z=z(x)$ and $z'=z(x')$ have the same active resource
rows and the same lower/free/upper box pattern. Take an independent
basis of their common active rows from $[C;I;-I]$. The row-space
solution $d$ of their active right-hand-side difference satisfies
\begin{equation}\label{eq:active-row-correction}
       \norm{d}_\infty\le D_b\Delta,\qquad
               e-d\text{ is tangent to all active rows},\quad e=z-z'.
\end{equation}
Indeed for the scaled integer basis $\mathsf A_J$ the solution is
$\mathsf A_J\trans(\mathsf A_J\mathsf A_J\trans)^{-1}$ times
the scaled right-hand-side difference. The latter has infinity norm
at most $d_0B_b\Delta$, and the cofactor bounds from
\cref{lem:resource-constants} give
$\norm{d}_\infty\le d_0N^2MV B_b\Delta\le D_b\Delta$.
Dependent active rows lie in the same row space, and their
right-hand-side differences are consistent because $e$ solves the complete
system.

By the polyhedral KKT conditions, both optimal gradients lie in the span
of the common active rows, so their difference is orthogonal to $e-d$.
Splitting the gradient difference through $\nabla F_{x'}(z')$ therefore gives
\[
 \begin{aligned}
 (\nabla F_x(z)-\nabla F_x(z'))\trans e
 &= (\nabla F_x(z)-\nabla F_{x'}(z'))\trans d\\
 &\quad +(\nabla F_{x'}(z')-\nabla F_x(z'))\trans e.
 \end{aligned}
\]
The first gradient difference has infinity norm at most
$A_z\norm{e}_\infty+A_x\Delta$, and the second at most $A_x\Delta$.
Adding the two Bregman inequalities at fixed leader $x$ now gives
\[
 2\mu\norm{e}_\infty^{P+1}
 \le(\nabla F_x(z)-\nabla F_x(z'))\trans e
 \le N(A_z\norm{e}_\infty+A_x\Delta)D_b\Delta
                                      +NA_x\Delta\norm{e}_\infty.
\]
If $\norm{e}_\infty\ge D_b\Delta$ and it is positive, division yields
\[
 \norm{e}_\infty^P\le\frac{N(A_zD_b+2A_x)}{2\mu}\Delta.
\]
Otherwise $\norm{e}_\infty\le D_b\Delta$. Since $x,x'\in X_F\subseteq
X\subseteq[0,1]^r$ (\cref{subsec:accuracy-model}), $\Delta\le1$, so
$D_b\Delta\le D_b\Delta^{1/P}$. Both cases therefore give the same-pattern
bound $\norm{e}_\infty\le C_0\Delta^{1/P}$ with rational constant
\[
 C_0=D_b+\max\{1,N(A_zD_b+2A_x)/(2\mu)\}.
\]
Continuity extends this bound to limits of any common-pattern subset.

It remains to control pattern changes along the segment
$x(t)=(1-t)x+tx'\in X_F$, $0\le t\le1$. The argument only bounds
their number; the algorithm does not construct their high-dimensional
description. A pattern has three choices per box coordinate and two per
resource row. For each pattern, its graph is described in $(t,z,\lambda)$
by box bound equalities or strict interior inequalities, active resource
equalities or strict slacks, $0\le\lambda\le\Lambda$, zero multipliers
on inactive resources, and the corresponding gradient signs or equations
for the box KKT conditions. These conditions are necessary and sufficient
by convexity, and use at most $3N+4k+2$ polynomial conditions of degree
at most $d=\max\{P,\deg\phi,2\}$. Real coefficients involving the two
fixed endpoint leaders do not affect the component count.

Convert each weak inequality $p\ge0$ to $p-u^2=0$, each strict
inequality $p>0$ to $pu^2-1=0$, and retain equations. Sum their squares.
The resulting real hypersurface has degree at most $D_*=2(d+2)$ in
at most $n_*=8(N+k+1)$ variables. Projection maps it onto the pattern
graph. The classical hypersurface component bound
\cite{Milnor1964} is at most $D_*(2D_*-1)^{n_*-1}$.
Continuous projection cannot increase the number of connected components,
so this also bounds the number of interval or point components of the
pattern's set of parameter values $t$. It applies to strict inequalities
because of the reciprocal-square variables, without a boundedness
assumption on those variables.

Consequently, the total component count over all patterns is at most the
explicit integer
\[
 \Gamma=3^N2^kD_*(2D_*-1)^{n_*}.
\]
The endpoints of those components partition $[0,1]$ into at most
$2\Gamma+1$ successive subintervals. On the interior of each choose a
pattern that occurs there; absence of component endpoints implies that
it occurs throughout that interior. Apply the same-pattern bound and
continuity at endpoints. Summing the bounds gives
\[
 \norm{z(x)-z(x')}_\infty
 \le C_0\sum_j[\Delta(t_j-t_{j-1})]^{1/P}
 \le C_0(2\Gamma+1)\Delta^{1/P}.
\]
Thus $C_z=C_0(2\Gamma+1)$ is a sufficient rational constant whose
logarithm is polynomial in the input and numerical degrees. Finally, the independent marginal $g(z)=z^P$ with
incentive $x$ has response $x^{1/P}$ near zero, so no larger uniform
H\"older exponent holds even for this one-variable subclass.
\end{proof}
